\documentclass[12pt, a4paper]{article}

% --- Paquetes --- % 
\usepackage[utf8]{inputenc}
\usepackage[T1]{fontenc}
\usepackage[spanish]{babel}
\usepackage{amsmath}
\usepackage{graphicx}
\usepackage{geometry}
\usepackage{listings}
\usepackage{hyperref}
\usepackage{xcolor}

% --- Configuración --- % 
\geometry{a4paper, margin=2.5cm}
\hypersetup{
    colorlinks=true,
    linkcolor=blue,
    filecolor=magenta,      
    urlcolor=cyan,
}

% Configuración para los listados de código Python
\definecolor{codegreen}{rgb}{0,0.6,0}
\definecolor{codegray}{rgb}{0.5,0.5,0.5}
\definecolor{codepurple}{rgb}{0.58,0,0.82}
\definecolor{backcolour}{rgb}{0.95,0.95,0.92}

\lstdefinestyle{mystyle}{
    backgroundcolor=\color{backcolour},
    commentstyle=\color{codegreen},
    keywordstyle=\color{magenta},
    numberstyle=\tiny\color{codegray},
    stringstyle=\color{codepurple},
    basicstyle=\ttfamily\footnotesize,
    breakatwhitespace=false,         
    breaklines=true,                 
    captionpos=b,                    
    keepspaces=true,                 
    numbers=left,                    
    numbersep=5pt,                  
    showspaces=false,                
    showstringspaces=false,
    showtabs=false,                  
    tabsize=2
}
\lstset{style=mystyle}

% --- Documento --- % 
\title{Análisis Bibliométrico de la Inteligencia Artificial Generativa\large \vspace{0.5em} 
    \normalsize Documento de Diseño y Arquitectura}
\author{[Tu Nombre]}
\date{\today}

\begin{document}

\maketitle

\begin{abstract}
    Este documento presenta el diseño, la arquitectura y los detalles de implementación de un sistema de software para el análisis bibliométrico en el dominio de la "inteligencia artificial generativa". El sistema automatiza la recolección y análisis de datos de publicaciones científicas, implementando algoritmos de similitud textual, análisis de frecuencia de palabras clave, clustering jerárquico y generación de visualizaciones. Adicionalmente, se documenta el rol de un asistente de Inteligencia Artificial (LLM) en el ciclo de desarrollo, detallando los prompts utilizados para la generación y refinamiento de código.
\end{abstract}

\tableofcontents
\newpage

\section{Introducción}
La bibliometría es una disciplina que permite analizar grandes volúmenes de datos derivados de la producción científica. Este proyecto tiene como propósito implementar un conjunto de herramientas de software que permitan el análisis bibliométrico y computacional sobre el dominio de la "inteligencia artificial generativa", a partir de las bases de datos disponibles en la Universidad del Quindío. El desarrollo se fundamenta en una serie de requerimientos funcionales que contemplan desde la limpieza de datos hasta la visualización y el análisis semántico de los mismos.

\section{Arquitectura del Proyecto}
El sistema se ha diseñado siguiendo una arquitectura modular, donde cada requerimiento funcional principal se encapsula en su propio script de Python. Esto facilita la mantenibilidad, la ejecución individual de cada tarea y la claridad del código.

\subsection{Flujo de Datos}
El flujo de datos del sistema es el siguiente:
\begin{enumerate}
    \item \textbf{Entrada de Datos:} Archivos crudos en formato `.bib` descargados de bases de datos científicas (ACM, ScienceDirect).
    \item \textbf{Procesamiento Inicial (Req. 1):} Un script inicial unifica los archivos, elimina duplicados basándose en el título y genera un corpus limpio (`articulos\_con\_titulo\_y\_abstract.bib`), que sirve como fuente única de verdad para los análisis posteriores.
    \item \textbf{Análisis Algorítmico (Req. 2-4):} Scripts dedicados (`requerimiento2.py`, `requerimiento3.py`, `requerimiento4.py`) consumen el corpus limpio para realizar análisis específicos.
    \item \textbf{Generación de Salidas (Req. 5):} Un script final (`requerimiento5.py`) genera todas las visualizaciones (imágenes `.png`, mapas interactivos `.html`) y las compila en un reporte final en formato `.pdf`.
\end{enumerate}

\subsection{Tecnologías y Librerías}
El proyecto se desarrolla enteramente en Python 3, utilizando un conjunto de librerías especializadas en ciencia de datos y procesamiento de lenguaje natural:
\begin{itemize}
    \item \textbf{Manipulación de Datos:} `bibtexparser`, `pandas`, `numpy`.
    \item \textbf{NLP y Machine Learning Clásico:} `scikit-learn`.
    \item \textbf{Modelos de IA:} `sentence-transformers` (para SBERT), `gensim` (para Doc2Vec).
    \item \textbf{Visualización:} `matplotlib`, `wordcloud`, `plotly`.
    \item \textbf{Reportes PDF:} `fpdf2`.
\end{itemize}

\section{Desarrollo de Requerimientos}

\subsection{Requerimiento 2: Similitud Textual}
\textbf{Objetivo:} Implementar seis algoritmos de similitud textual (cuatro clásicos y dos con IA) para comparar los abstracts de los artículos.

\textbf{Implementación:} Se creó el script `requerimiento2.py`. Este script carga los datos y define una función para cada uno de los siguientes algoritmos, incluyendo en los comentarios del código una explicación detallada de su funcionamiento:
\begin{itemize}
    \item Distancia de Levenshtein
    \item Similitud de Jaccard
    \item TF-IDF con Similitud del Coseno
    \item Bag of Words (BoW) con Similitud del Coseno
    \item Sentence-BERT (SBERT)
    \item Doc2Vec
\end{itemize}
El script principal compara los dos primeros artículos del corpus como demostración.

\subsection{Requerimiento 3: Frecuencia de Palabras Clave}
\textbf{Objetivo:} Calcular la frecuencia de un listado de palabras clave predefinidas y, adicionalmente, generar un nuevo listado de términos relevantes a partir de los abstracts.

\textbf{Implementación:} El script `requerimiento3.py` realiza dos tareas:
\begin{enumerate}
    \item \textbf{Frecuencia de Palabras Dadas:} Utiliza expresiones regulares (`re.findall`) para contar las apariciones de cada término de la lista predefinida en la totalidad de los abstracts.
    \item \textbf{Generación de Nuevas Palabras:} Aplica la técnica TF-IDF sobre todo el corpus para identificar los 15 términos (unigramas y bigramas) con mayor relevancia. Finalmente, compara esta lista generada con la original para ofrecer una métrica de "precisión".
\end{enumerate}

\subsection{Requerimiento 4: Clustering Jerárquico}
\textbf{Objetivo:} Implementar tres algoritmos de agrupamiento jerárquico para construir dendrogramas que representen la similitud entre los abstracts.

\textbf{Implementación:} El script `requerimiento4.py`:
\begin{enumerate}
    \item \textbf{Vectorización:} Transforma cada abstract en un vector semántico utilizando un modelo Sentence-BERT pre-entrenado (`all-MiniLM-L6-v2`).
    \item \textbf{Clustering:} Aplica los algoritmos de clustering jerárquico utilizando `scipy.cluster.hierarchy.linkage` con tres métodos diferentes: `ward`, `complete` y `average`.
    \item \textbf{Visualización:} Genera un dendrograma para cada método utilizando `matplotlib` y lo guarda como un archivo `.png`.
\end{enumerate}

\subsection{Requerimiento 5: Análisis Visual}
\textbf{Objetivo:} Crear un conjunto de visualizaciones y exportarlas a un reporte en PDF.

\textbf{Implementación:} El script `requerimiento5.py` centraliza la generación de todos los artefactos visuales:
\begin{itemize}
    \item \textbf{Nube de Palabras:} Generada con la librería `wordcloud`.
    \item \textbf{Líneas de Tiempo:} Gráficos de barras creados con `pandas` y `matplotlib`.
    \item \textbf{Mapa de Calor:} Un mapa coroplético interactivo generado con `plotly`, que requiere la extracción y estandarización de los nombres de países desde los datos de autor.
    \item \textbf{Reporte PDF:} Utiliza `fpdf2` para compilar las imágenes `.png` generadas y una referencia al mapa `.html` en un documento `reporte\_visual.pdf`.
\end{itemize}

\section{Fundamentos y Uso de IA en el Desarrollo}
Para acelerar el ciclo de desarrollo, generar código base, depurar errores y obtener explicaciones algorítmicas, se utilizó un Asistente de IA (Large Language Model). A continuación, se presentan ejemplos de los prompts formales y estructurados que guiaron la generación de código para cada requerimiento. Estos prompts fueron diseñados para ser claros, específicos y proporcionar todo el contexto necesario para que la IA generara una respuesta útil y precisa.

\subsection{Prompts para el Requerimiento 2}
\begin{lstlisting}[language=bash, caption={Prompt para el algoritmo TF-IDF}]
# Instrucción para el Asistente de IA

Contexto: Estoy trabajando en el script `requerimiento2.py`, que ya implementa la Distancia de Levenshtein y la Similitud de Jaccard. El script carga artículos desde un archivo .bib.

Tarea: Integra el algoritmo de "TF-IDF con Similitud del Coseno".

Requisitos:
1.  Añade las importaciones necesarias desde `sklearn`.
2.  Escribe una función `calcular_similitud_tfidf_coseno` que reciba dos textos y el corpus completo de abstracts (necesario para el cálculo de IDF).
3.  Dentro de la función, inicializa un `TfidfVectorizer`, ajústalo al corpus y luego transforma los dos textos de entrada.
4.  Calcula la similitud del coseno entre los dos vectores resultantes.
5.  Añade un bloque de comentarios explicando en detalle el funcionamiento de TF-IDF y la Similitud del Coseno.
6.  En el bloque principal, llama a esta nueva función y muestra el resultado.
\end{lstlisting}

\subsection{Prompts para el Requerimiento 3}
\begin{lstlisting}[language=bash, caption={Prompt para la generación de nuevas palabras clave}]
# Instrucción para el Asistente de IA

Contexto: El script `requerimiento3.py` ya calcula la frecuencia de un listado predefinido de palabras.

Tarea: Añade una nueva funcionalidad para generar un listado de 15 palabras clave nuevas directamente desde los abstracts del corpus.

Requisitos:
1.  Crea una función `generar_nuevas_palabras_tfidf` que reciba el corpus de abstracts.
2.  Dentro de la función, utiliza `TfidfVectorizer` de `sklearn` para procesar el texto. Configúralo para usar `stop_words='english'` y un `ngram_range=(1, 2)` para capturar unigramas y bigramas.
3.  Calcula la suma de los puntajes TF-IDF para cada término a través de todos los documentos.
4.  Identifica los 15 términos con el puntaje más alto.
5.  Devuelve esta lista de 15 términos.
6.  En el bloque principal, llama a esta función y presenta los resultados en la consola.
\end{lstlisting}

\subsection{Prompts para el Requerimiento 4}
\begin{lstlisting}[language=bash, caption={Prompt para el clustering jerárquico}]
# Instrucción para el Asistente de IA

Contexto: Se necesita crear un nuevo script `requerimiento4.py` para el clustering de artículos.

Tarea: Implementa una solución completa para el clustering jerárquico de abstracts.

Requisitos:
1.  Carga los artículos desde el archivo .bib.
2.  Utiliza el modelo `all-MiniLM-L6-v2` de la librería `sentence-transformers` para convertir cada abstract en un embedding vectorial.
3.  Crea una función que reciba los embeddings, las etiquetas (títulos de artículos) y un método de clustering (`ward`, `complete`, `average`).
4.  Dentro de la función, utiliza `scipy.cluster.hierarchy.linkage` para realizar el clustering.
5.  Posteriormente, utiliza `matplotlib` y `scipy.cluster.hierarchy.dendrogram` para generar un dendrograma y guardarlo como un archivo PNG.
6.  En el bloque principal, itera sobre los tres métodos de clustering y genera un dendrograma para cada uno.
\end{lstlisting}

\subsection{Prompts para el Requerimiento 5}
\begin{lstlisting}[language=bash, caption={Prompt para el mapa de calor geográfico}]
# Instrucción para el Asistente de IA

Contexto: El script `requerimiento5.py` ya genera una nube de palabras y líneas de tiempo.

Tarea: Añade la funcionalidad para generar un mapa de calor geográfico interactivo.

Requisitos:
1.  Importa las librerías `plotly.express` y `pycountry_convert`.
2.  Crea una función `generar_mapa_calor`.
3.  Dentro de la función, extrae el país de la afiliación del primer autor. Dado que el campo `author` es texto libre, implementa una heurística que tome la última parte de la cadena de texto y intente convertirla a un nombre de país estandarizado.
4.  Cuenta el número de publicaciones por país.
5.  Convierte los nombres de los países a códigos ISO ALPHA-3, requeridos por Plotly.
6.  Utiliza `px.choropleth` para generar un mapa del mundo, coloreando cada país según su número de publicaciones.
7.  Guarda el mapa interactivo como un archivo HTML.
\end{lstlisting}

\section{Conclusión}
Este proyecto ha logrado implementar con éxito un pipeline completo para el análisis bibliométrico. Se ha demostrado la viabilidad de combinar técnicas clásicas de procesamiento de texto con modelos modernos de IA para extraer insights significativos de la literatura científica. La arquitectura modular del proyecto y la documentación exhaustiva, incluyendo la formalización del proceso de desarrollo asistido por IA, aseguran su reproducibilidad y extensibilidad futura.

\end{document}
